\subsection{FnType and Asts List}
\par The other place where GADTs and type families are needed are for applying constructors to core expression arguments.
In order for the typechecker to allow this, it needs to know the type of the constructor 
(i.e what arguments it needs for construction and what the result of the construction is).
Since the constructor could take an arbitrary number of arguments, 
its type had to be computed with a type family (see figure~\ref{fig:fntype}).
The FnType is parametrized with $\mathit{args}$ and $r$ where $\mathit{args}$ is a list of $n$ and $a$ pairs describing
the depths and root element type of each argument in order, and $r$ is result type.
For example, the list cons constructor has type \code{FnType [(0, Int), (1, Int)] (NestedList 1 Int)}. 
To get its actual type, the typechecker applies the FnType rules to get \code{(NestedList 0 Int) -> (NestedList 1 Int) -> (NestedList 1 Int)}.
Then it applies the NestedList rules to get the final computed type of \code{Int -> [Int] -> [Int]}.

This type family is used in conjunction with the Asts List GADT (see figure~\ref{fig:astslist}) to fully realize this functionality.
The Asts list encodes the depth an root element types of each of its elements in its own type as the type argument $\mathit{args}$.
This means that the the constructor for data can look like this:
\codeln{Data :: FnType args a -> Asts args -> Ast 0 a}
Since the $\mathit{args}$ from FnType and the Asts list are the same, the typechecker knows that when the elements of the list
are used to generate values, those values can be passed to the function in order.

\begin{figure}
    \begin{hskcode}
    type family FnType (args :: [(Nat, Type)]) r where
        FnType '[] r = r
        FnType ('(n, a) ': xs) r = NestedList n a -> FnType xs r
    \end{hskcode}
    \caption{Haskell code for the FnType}
    \label{fig:fntype}
\end{figure}
\begin{figure}
    \begin{hskcode}
    data Asts (args :: [(Nat, Type)]) where
        AstsN     :: Asts '[]
        (::>)     :: Ast n a -> Asts xs -> Asts ('(n, a) ': xs)
    \end{hskcode}
    \caption{Haskell code for the NestedList type family}
    \label{fig:nestedlist}
\end{figure}
